\documentclass[12pt]{scrartcl}
\addtokomafont{disposition}{\boldmath}
% \documentclass[12pt]{article}
% \usepackage[v1]{subfiles} %%

\usepackage[utf8]{inputenc}
% \usepackage[croatian]{babel}
\usepackage{datetime}
\usepackage[a4paper,margin=0.75in]{geometry}
\setlength\baselineskip{15pt}

\usepackage{amsfonts}
\usepackage{amsmath}
\usepackage{amssymb}
\usepackage{amsthm}
\usepackage{csquotes}
\usepackage{tikz}
\usepackage{arydshln}
\pagenumbering{gobble}
\usepackage{float}
\usepackage[dvipsnames]{xcolor}
%\usepackage{mathptmx} % times new roman
% \usepackage{mathpazo} % tex gyre pagella
% \usepackage{baskervaldx}
% \usepackage[baskervaldx]{newtxmath}
% \usepackage[T1]{fontenc}
% \fontseries{ub}\selectfont

\usepackage{breqn}
\usepackage{thmtools}
\usepackage{multirow}
\usepackage{mathtools}
\usepackage[unicode]{hyperref}

\usepackage{booktabs}
\usepackage{listings}

%%%

% Save original macros
% --------------------
\usepackage{letltxmacro}

\LetLtxMacro\OriginalLongrightarrow\Longrightarrow
\LetLtxMacro\OriginalLongleftarrow\Longleftarrow

% Implement new macros
% --------------------
\usepackage{trimclip}
\DeclareRobustCommand\Longrightarrow{\NewRelbar\joinrel\Rightarrow}
\DeclareRobustCommand\Longleftarrow{\Leftarrow\joinrel\NewRelbar}

\makeatletter
\DeclareRobustCommand\NewRelbar{%
  \mathrel{%
    \mathpalette\@NewRelbar{}%
  }%
}
\newcommand*\@NewRelbar[2]{%
  % #1: math style
  % #2: unused
  \sbox0{$#1=$}%
  \sbox2{$#1\Rightarrow\m@th$}%
  \sbox4{$#1\Leftarrow\m@th$}%
  \clipbox{0pt 0pt \dimexpr(\wd2-.6\wd0) 0pt}{\copy2}%
  \kern-.2\wd0 %
  \clipbox{\dimexpr(\wd4-.6\wd0) 0pt 0pt 0pt}{\copy4}%
}
\makeatother
%%%


\definecolor{seagreen}{HTML}{21b2aa}
\definecolor{magenta}{HTML}{b2217f}
\definecolor{gold}{HTML}{ca9520}
\definecolor{red}{HTML}{da272f}
\definecolor{blue}{HTML}{4682b4}
\definecolor{navy}{HTML}{7b000b} %% orig 06038d % onda 7b0000 % onda 7b000b
\definecolor{nred}{HTML}{650015}
\definecolor{ngreen}{HTML}{156515}

%% 
\definecolor{highlightcolor}{named}{ngreen}
\addtokomafont{disposition}{\color{highlightcolor}}
\renewcommand{\emph}[1]{\textbf{\textcolor{highlightcolor}{#1}}}
%%

\definecolor{codebg}{HTML}{f0f0f0}
\definecolor{commentfg}{HTML}{333333}

\lstdefinestyle{mystyle}{
	backgroundcolor=\color{codebg},
	commentstyle=\color{commentfg},
	keywordstyle=\color{navy},
	stringstyle=\color{commentfg},
	basicstyle=\ttfamily\footnotesize,
	breakatwhitespace=false,
	breaklines=true,
	captionpos=b,
	keepspaces=true,
	numbers=none,
	showspaces=false,
	showstringspaces=false,
	showtabs=false,
	tabsize=2
}
\lstset{style=mystyle}

\hypersetup{
    colorlinks,
    linkcolor=highlightcolor,
    citecolor=highlightcolor,
    urlcolor=highlightcolor
}

\usepackage{enumitem}
\makeatletter
\providecommand*{\blx@noerroretextools}{}
\makeatother

\usepackage[
%	style=alphabetic,
	style=numeric,
	giveninits
	% firstinits=true,
	% terseinits=true,
	% giveninits=true
]{biblatex}
%\addbibresource{literatura.bib}
\let\oldcite\cite
\renewcommand{\cite}[2][]{%
    \def\tmp{#1}%
    \ifx\tmp\empty
        \mbox{\oldcite{#2}}%
    \else
        \mbox{\oldcite[\tmp]{#2}}%
    \fi
}
\usepackage{pgfplots}
\pgfplotsset{compat=1.15}
\usepackage{mathrsfs}
\usetikzlibrary{arrows}
	
\usepackage[section]{placeins}
\declaretheorem{teorem}
% \numberwithin{teorem}{subsection}
% \numberwithin{equation}{section}
% \numberwithin{figure}{section}
% \numberwithin{table}{section}
\declaretheorem[sibling=teorem, style=plain]{lemma}
\declaretheorem[style=remark, sibling=teorem]{note}
%\declaretheorem[style=remark, sibling=teorem]{komentar*}
\declaretheorem[style=definition, sibling=teorem]{definition}
\declaretheorem[sibling=teorem]{corollary}
% \declaretheorem[style=remark, sibling=teorem]{zadatak}
\declaretheorem[sibling=teorem]{proposition}
\declaretheorem[sibling=teorem]{example}
\declaretheorem[sibling=teorem]{problem}
\declaretheorem[sibling=teorem]{exercise}

\newcommand{\ub}{\rightarrow \infty}
\newcommand{\un}{\rightarrow 0}
% \newcommand{\norm}[1]{\left\lVert#1\right\rVert} 
% \newcommand{\given}{\, \middle\vert \,}
\newcommand{\given}{\, \vert \,}
\newcommand{\st}{\, \colon \,}
\newcommand{\D}{\,\mathrm d}
\newcommand{\ani}{{^\perp}}
\newcommand{\ran}{\mathrm{ran}\,}
\newcommand{\bop}{\mathrm{B}}
% \newcommand{\skp}[2]{\left\langle {#1}, {#2}  \right\rangle}
\newcommand{\jpod}{\stackrel{\mathrm d}{=}}
\newcommand{\jgs}{\stackrel{\mathrm{g.s.}}{=\joinrel=}}
\newcommand{\jas}{\stackrel{\mathrm{a.s.}}{=\joinrel=}}
%\newcommand{\jzbog}[1]{\stackrel{\eqref{#1}}{=}}
\newcommand{\wkonv}{\Rightarrow}
% \newcommand{\stackrelgd}[2]{\stackrel{\stackrel{ {#1}} {\longrightarrow}} { {\scriptstyle #2} } }
\newcommand{\stackrell}[2]{\stackrel{\scriptstyle {#1}}{#2}}
\newcommand{\konvp}{\stackrel{\mathbb P}{\longrightarrow}}
\newcommand{\konvw}{\stackrel{\mathrm w}{\longrightarrow}}
\newcommand{\konvd}{\stackrel{\mathrm d}{\longrightarrow}}
\newcommand{\konvgs}{\stackrel{\mathrm{g.s.}}{\longrightarrow}}
\newcommand{\konvas}{\stackrel{\mathrm{a.s.}}{\longrightarrow}}
\newcommand{\konvv}{\stackrel{\mathrm{v}}\longrightarrow}
\newcommand{\konvL}[1]{\stackrel{\L^{#1} }\longrightarrow}
\DeclareMathOperator{\card}{card}
\DeclareMathOperator{\var}{Var}
\DeclareMathOperator{\supp}{supp}
\DeclareMathOperator{\im}{im}
\DeclareMathOperator{\olspan}{\overline{span}}
\DeclareMathOperator{\tg}{tg}
\DeclareMathOperator{\ctg}{ctg}
\DeclareMathOperator{\arctg}{arctg}
\DeclareMathOperator{\Arsh}{Arsh}
\DeclareMathOperator{\discont}{discont}
\DeclareMathOperator{\Int}{Int}
\DeclareMathOperator{\cl}{Cl}
\DeclareMathOperator{\dist}{d}
\DeclareMathOperator{\re}{Re}
\let\im\relax
\DeclareMathOperator{\im}{Im}
\DeclareMathOperator{\sgn}{sgn}
% \DeclareMathOperator{\ln}{\ln}

\newcommand{\ol}[1]{\overline{#1}}
\newcommand{\wh}[1]{\widehat{#1}}
\newcommand{\wt}[1]{\widetilde{#1}}
\renewcommand{\th}[1]{\hat{#1}} % drukciji hat za tekst?
% \newcommand{\eqdef}{\ {=\colon} \ } 
% \newcommand{\defeq}{\ {\colon\!=} \ }

\newcommand{\R}{\mathbb{R}}
\newcommand{\N}{\mathbb{N}}
\newcommand{\Z}{\mathbb{Z}}
\newcommand{\Q}{\mathbb{Q}}
\newcommand{\C}{\mathbb{C}}
\renewcommand{\P}{\mathbb{P}}
\newcommand{\Ptrans}{\P}  % trazicijska funkcija Markovljevog procesa
\newcommand{\E}{\mathbb{E}}
\def\B{\mathrm{B}}
\def\K{\mathrm{K}}
\def\M{\mathrm{M}}
\def\CC{\mathrm{C}}
\def\RR{\mathrm{R}}
\def\H{\mathcal{H}}
\def\K{\mathcal{K}}
% \def\Z{\mathcal{Z}}
\def\S{\mathcal{S}}
\def\T{\mathcal{T}}
\def\U{\mathcal{U}}

\newcommand{\cz}{Calder\' on--Zygmund}
\newcommand{\holder}{H\" older}
\newcommand{\cadlag}{c\` adl\` ag}
\newcommand{\levy}{L\' evy}
\newcommand{\ito}{It\= o}
\def\L{\mathrm{L}}
\def\w{\mathrm{w}}

%%
% \newcommand{\oo}[1]{\left\langle {#1} \right\rangle}
% \newcommand{\oc}[1]{\left\langle {#1} \right]}
% \newcommand{\co}[1]{\left[ {#1} \right\rangle}
% \newcommand{\cc}[1]{\left[ {#1} \right]}

% \newcommand{\ob}[1]{\left( {#1} \right)}
% \newcommand{\ug}[1]{\left[ {#1} \right]}
% \newcommand{\vit}[1]{\left\{ {#1} \right\}}
\DeclarePairedDelimiter{\vit}{\{}{\}}
\let\oldvit\vit
\renewcommand{\vit}[1]{\oldvit*{#1}}

\DeclarePairedDelimiter{\ug}{[}{]}
\let\oldug\ug
\renewcommand{\ug}[1]{\oldug*{#1}}

\DeclarePairedDelimiter{\ob}{(}{)}
\let\oldob\ob
\renewcommand{\ob}[1]{\oldob*{#1}}

\DeclarePairedDelimiter{\abs}{\lvert}{\rvert}
\let\oldabs\abs
\renewcommand{\abs}[1]{\oldabs*{#1}}

\DeclarePairedDelimiter{\norm}{\lVert}{\rVert}
\let\oldnorm\norm
\renewcommand{\norm}[1]{\oldnorm*{#1}}
\newcommand{\dnorm}{\norm{\, \cdot \,}}

\DeclarePairedDelimiter{\skp}{\langle}{\rangle}
\let\oldskp\skp
\renewcommand{\skp}[1]{\oldskp*{#1}}

\DeclarePairedDelimiter{\oo}{\langle}{\rangle}
\let\oldoo\oo
\renewcommand{\oo}[1]{\oldoo*{#1}}

\DeclarePairedDelimiter{\oc}{\langle}{]}
\let\oldoc\oc
\renewcommand{\oc}[1]{\oldoc*{#1}}

\DeclarePairedDelimiter{\co}{[}{\rangle}
\let\oldco\co
\renewcommand{\co}[1]{\oldco*{#1}}

\DeclarePairedDelimiter{\cc}{[}{]}
\let\oldcc\cc
\renewcommand{\cc}[1]{\oldcc*{#1}}


\newcommand{\up}[1]{\mathrm{#1}}
\renewcommand{\cal}[1]{\mathcal{#1}}

\newcommand{\pd}{\subseteq}
\newcommand{\nd}{\supseteq}

\newcommand{\rz}{\setminus}

\renewcommand{\le}{\leqslant}
\renewcommand{\ge}{\geqslant}

\newcommand{\ime}[2]{\, {#1}({\mathrm d {#2}})}

\newcommand{\io}{\ \text{i.o.}}
\newcommand{\ult}{\ \text{ult.}}

\newcommand{\whi}{; \,}
\newcommand{\zerone}{\( 0 \)-\( 1 \)}
\newcommand{\mii}{_{-\infty}^\infty}
\newcommand{\nderi}[1]{^{(#1)}}

\newcommand{\implybox}[2]{\fbox{\emph{ ({#1}) \( \implies \) ({#2}) }}}

\newcommand\restr[2]{{% we make the whole thing an ordinary symbol
  \left.\kern-\nulldelimiterspace % automatically resize the bar with \right
  #1 % the function
  \vphantom{\big|} % pretend it's a little taller at normal size
  \right|_{#2} % this is the delimiter
  }}

\newcommand\neza{\protect\mathpalette{\protect\independenT}{\perp}}
\def\independenT#1#2{\mathrel{\rlap{$#1#2$}\mkern2mu{#1#2}}}

%%
%
\makeatletter
\DeclareRobustCommand{\lasymp}{\lg@asymp{<}}
\DeclareRobustCommand{\gasymp}{\lg@asymp{>}}

\newcommand{\under@asymp}[1]{\clipbox{0pt 0pt 0pt {0.5\height}}{$\m@th#1\asymp$}}
\newcommand{\lg@asymp}[1]{\mathrel{\mathpalette\lg@asymp@{#1}}}
\newcommand{\lg@asymp@}[2]{%
  \vcenter{%
    \offinterlineskip
    \m@th
    \ialign{%
      \hfil##\hfil\cr
      $#1#2$\cr
      \under@asymp{#1}\cr
    }%
  }%
}
\makeatother

%% specificno za fajl
\newcommand{\sal}{\( \sigma \)-algebra}
%%


%\newcommand{\distr}[3]{\left\{ x \in {#1} \st \left| {#2}(x) \right| \ge {#3} \right\}}
% \newcommand{\abs}[1]{\left| {#1} \right|}
% \newcommand{\fave}[2]{\left\langle {#1} \right\rangle_{#2}}
\let\etoolboxforlistloop\forlistloop % save the good meaning of \forlistloop
\usepackage{autonum}
\let\forlistloop\etoolboxforlistloop % restore the good meaning of \forlistloopi

\newcommand{\asteriskline}{
\vskip20pt
\centerline{
\Large{
\textasteriskcentered\hspace{.1em}
\raisebox{.87ex}\textasteriskcentered\hspace{.1em}
\textasteriskcentered
}}
\vskip20pt
}

\newcommand{\astfootnote}[1]{%
\let\oldthefootnote=\thefootnote%
\setcounter{footnote}{0}%
\renewcommand{\thefootnote}{\fnsymbol{footnote}}%
\footnote{#1}%
\let\thefootnote=\oldthefootnote%
}

\begin{document}
\pagenumbering{gobble}
\title{Notes on two walks with noncollinear drifts}
% \author{Luka Šimek}
% \date{\today}

%\date{\today}
\maketitle
% \vspace{-5em}
% \begin{flushleft}
% 	These materials...
% \end{flushleft}
% \vspace{-.5em}


% \vspace{.25em}  % Add some space before signature?
%\tableofcontents
%\newpage
\pagenumbering{arabic}

We use the notation:
\begin{itemize}
	\item polar form of the walk: \( S_k\nderi j = \norm {S_k \nderi j} e_{ \varphi_k\nderi j  } \); then
	      \begin{equation}
		      \skp{S_k \nderi j, e_\theta} = \norm{S_k \nderi j} \cos \ob{\varphi_k \nderi j - \theta}
	      \end{equation}
	\item for the index \( k \) where the max of \( \skp {S_k\nderi j, e_\theta} \) is achieved, we write \( k^*(\theta, n, j) \coloneq k \)
	      % \item \( Q_j \) for quadrants of \( \R^2 \).
\end{itemize}

\section{Lemmas}
By the SLLN, all points \( S_k\nderi 1 \) for \( k \ge n \) where \( n=n(\omega, \delta) \) is a random and a.s.\ finite \( \in \N \)
are contained in the cone
\begin{equation}\label{eq:cone1}
	\up{Co}\nderi1 = \vit{ r e_\theta \st r \ge 0, \ \theta \in \oo{\frac\pi4 - \delta, \frac\pi4 + \delta} }
\end{equation}
and likewise for the second walk.

\begin{lemma}\label{lm:21}
	A.s.\ for large enough \( n \)
	it holds that
	\begin{equation}
		S_{k^*(\theta, n, 1)}\nderi 1 \in \up{Co}\nderi 1, \quad \forall
		\theta \in \oo{-\frac\pi4+2\delta, \frac{3\pi}4-2\delta  },
	\end{equation}

\end{lemma}

\begin{proof}
	Write
	\begin{equation}\label{eq:max12}
		\begin{gathered}
			\max_{1 \le m \le n} \skp{S_m \nderi 1, e_\theta} =
			\max \vit{\max_{S_m \nderi 1 \not \in \up{Co} \nderi 1} \skp{S_m \nderi 1, e_\theta}, \max_{S_m \nderi 1\in \up{Co} \nderi 1} \skp{S_m \nderi 1, e_\theta} }
		\end{gathered}
	\end{equation}
	The first part of~\eqref{eq:max12} will a.s.\ at some point remain constant, and can therefore be bounded by the maximum of norms (making it independent of \( \theta \) as well as \( n \)).

	For the second part note first
	\begin{equation}
		\skp {S_m \nderi 1, e_\theta } = \norm{S_m \nderi 1} \cos \ob{\theta - \varphi_m \nderi 1}.
	\end{equation}
	When \( S_m \nderi 1 \in \up{Co}\nderi 1 \) we have \( \cos \ob{\theta - \varphi_m \nderi 1} > 0 \) and \( \norm{S_m \nderi 1} \konvas \infty \)
	meaning that the maximum in~\eqref{eq:max12} is a.s.\ always generated by the second part for large \( n \).

	Moreover,
	\begin{equation}\label{eq:strposcon}
		\cos \ob{\theta - \varphi_m \nderi 1} > \cos \ob{-\frac \pi 2 + \delta}
	\end{equation}
	which is a strictly positive constant, meaning that we can choose an \( n \) which
	is \enquote{good} for all such \( \theta \) simultaneously.
\end{proof}

\begin{note}\label{n1}
	% s manjim intervalom
	When
	\begin{equation}
		\theta \in \oo{-\frac \pi 4 + \delta, \frac \pi 4 - \delta },
	\end{equation}
	we cannot repeat the same argument, as the RHS in~\eqref{eq:strposcon} is \( 0 \).
\end{note}

\begin{note}\label{n2}
	For any fixed \( \theta \in \oo{-\pi/4, 3\pi/4} \) we can retroactively choose a small enough \( \delta \) to claim \( S_{k^*(\theta,n,1)} \in \up{Co}\nderi 1  \) a.s.\ for large enough \( n \), but we cannot do so for all \( \theta \) simultaneously.
\end{note}

\begin{note}\label{n3}
	Whenever \( \theta \) is instead such that \( \cos \ob{\theta - \varphi_m \nderi 1} < 0 \) inside the cone, the maximum
	is generated either by \( 0 \) or by a point outside the cone, meaning that the maximum a.s.\ eventually \enquote{freezes}. To make the analogous claim to lemma~\ref{lm:21}
	% the simpler
	\begin{equation}
		\theta \in \oo{ \frac{3\pi}4 + \delta , \frac{7\pi}4 -\delta \equiv -\frac \pi4-\delta}
	\end{equation}
	is enough, since the \( \skp{S_n \nderi 1, e_\theta} \) become negative for large enough \( n \), and
	the magnitude of the cosine (and therefore \( \theta \) itself) plays no part here.
\end{note}

\section{Main part}
Consider the limiting behaviour of the integral
\begin{equation}\label{eq:intmaincijeli}
	\frac 1{n^{3/2}} \int_0^{\pi} \ug{M_{n,2}(\theta)^2 - M_{n,1}(\theta)^2 } \cdot 1_{\vit{M_{n,1}(\theta) < M_{n,2}(\theta)}} \D \theta.
\end{equation}

\emph{For \( \theta \in \cc{\varepsilon, \pi/4 + \delta} \).} We can bound the indicator uniformly in \( \theta \) using
\begin{equation}
	\vit{M_{n,1}(\theta) < M_{n,2}(\theta)} \pd
	\vit{\min_{\theta \in \cc{\varepsilon, \pi/4 + \delta}} M_{n,1}(\theta)
		< \max_{\theta \in \cc{\varepsilon, \pi/4 + \delta}} M_{n,2}(\theta)}, \quad \forall \theta \in \cc{\varepsilon, \pi/4 + \delta}.
\end{equation}
For \( n \) large enough in the sense of lemma~\ref{lm:21} for both walks we claim:
\begin{equation}
	\min_{\theta \in \cc{\varepsilon, \pi/4+\delta}} M_{n,1}(\theta) = M_{n,1}(\varepsilon), \quad
	\max_{\theta \in \cc{\varepsilon, \pi/4+\delta}} M_{n,2}(\theta) = M_{n,2}(\varepsilon).
\end{equation}

We give the proof for the first walk. Assuming the contrary, there is a \( \theta > \varepsilon \)
and indices \( j,k \) generating maxima for \( \theta, \varepsilon \) respectively such that
\begin{equation}
	\skp{S_j \nderi 1, e_\theta} < \skp{S_k \nderi 1, e_\varepsilon}.
\end{equation}
But then
\begin{equation}
	\skp {S_k \nderi 1,  e_\theta} \le \skp{S_j \nderi 1, e_\theta} < \skp{S_k \nderi 1, e_\varepsilon}
\end{equation}
implying
\begin{equation}
	\cos \ob{\varphi_k \nderi 1 - \theta} < \cos \ob{\varphi_k \nderi 1 - \varepsilon}
\end{equation}
yielding a contradiction since \( \varphi_k \nderi 1 \) corresponds to the cone \( \up{Co} \nderi1 \).

Since by the SLLN
\begin{equation}
	1_{\vit{M_{n,1}(\varepsilon) < M_{n,2}(\varepsilon)}} \konvas 0, \quad \forall \varepsilon > 0,
\end{equation}
this proves convergence a.s.\ to \( 0 \).
The proof for the second walk is analogous.\footnote{
	note that to apply lemma~\ref{lm:21} we'd need
	the upper endpoint to be \( \pi/4-2\delta \) instead. However, for this proof
	to work we only need the result at the fixed \( \theta = \varepsilon \). The upper
	endpoint \( \pi/4+\delta \) is chosen for the sake of the next part.}

\emph{For \( \theta \in \cc{0,\varepsilon} \)}. Note that above we've also shown that
\begin{equation}
	\theta \mapsto M_{n,1}(\theta)
\end{equation}
is increasing on the domain \( \cc{0, \varepsilon} \) (a.s.\ for a large-enough random \( n \)).
We prove \( \L^1 \)-convergence.

\emph{For \( \theta \in \oc{\pi/4+\delta, \pi} \)} Here \( M_{n,2}(\theta) \) will a.s.\ freeze (at the same time for all \( \theta \)), so simply
\begin{equation}
	M_{n,2}(\theta)^2  - M_{n,1}(\theta)^2
	\le M_{n,2}(\theta)^2 \le \max_{S_m \nderi 2 \not \in \up{Co} \nderi 2} \norm{S_m \nderi 2}^2,
	\quad \text{on}\quad \vit{M_{n,1}(\theta) < M_{n,2}(\theta)}
\end{equation}
makes a.s.\ convergence to \( 0 \) under \( n^{-3/2} \) scaling obvious.

\emph{For \( \theta \in \oc{-\pi,0} \)}, the situation is symmetrical to the previous 3 points.

\section{Derivatives}
Writing \( S_n \) for a generic walk, Pythagoras' theorem gives
\begin{equation}
	\norm {S_{k^*(\theta, n)}}^2 = M_n(\theta) ^2 + M_n'(\theta)^2,
\end{equation}
meaning that
\begin{equation}
	M'_{n,2}(\theta)^2 - M'_{n,1}(\theta)^2 =
	\ug{ \norm{S_{k^*(\theta,n,2)} \nderi 2}^2 - \norm{S_{k^*(\theta,n,1)} \nderi 1} ^2 } -
	\ug{M_{n,2}(\theta)^2 - M_{n,1}(\theta)^2}
\end{equation}

The norms can be bounded using \( M_n \). First note
\begin{equation}
	\norm { S_{k^*(\theta, n, 1)}\nderi 1 } =
	\frac{M_{n, 1}(\theta)}{\cos \ob{\theta - \varphi_{k^*(\theta, n, 1)}\nderi1 }}
\end{equation}
For example when
\( \theta \in \oo{0,\varepsilon} \) we can guarantee
\begin{equation}
	\varphi_{k^*(\theta, n, 1)} \nderi1  -\theta \in \oo{\frac\pi4 - \delta - \varepsilon, \frac\pi4 + \delta}
\end{equation}
(again using lemma~\ref{lm:21}). Doing the same for the second walk we conclude
\begin{equation}
	\norm{S_{k^*(\theta,n,2)}\nderi2} +
	\norm{S_{k^*(\theta,n,1)}\nderi1}
	\lasymp_{\delta, \varepsilon}
	M_{n,2}(\theta) + M_{n,1}(\theta).
\end{equation}
Thus when bounding the sum, the previous results are sufficient from here.

We cannot do the same for the difference
\begin{equation}
	\norm{S_{k^*(\theta,n,2)}\nderi2} -
	\norm{S_{k^*(\theta,n,1)}\nderi1}
\end{equation}
as the bounds will then yield, for positive real numbers \( a > b \)
\begin{equation}
	a \norm2 - b\norm1 = b \ob{\norm 1-\norm 2} + (a-b)\norm 2
\end{equation}
which is \( O(n) \).

Instead,
\newcommand{\Sj}[1]{S_{k^*(\theta,n,#1)}\nderi#1}
\begin{align}
	\norm{\Sj2} - \norm{\Sj1} =
	\ob{\norm{S_n \nderi 2} - \norm{S_n \nderi 1} }
	+ \ob{\norm{\Sj2} - \norm{S_n \nderi 2}} -
	\ob{\norm{\Sj1} - \norm{S_n \nderi 1}}
\end{align}
Intuitively,
\begin{itemize}
	\item the 1st part will be \( O(\sqrt n) \) as the drifts have the same magnitudes
	      and these cancel each other out;
	\item the 2nd and 3rd part don't grow with \( n \), they should be
	      \( O(1) \).
\end{itemize}

If there is a need to remove dependence on \( \theta \), consider the
following. We'd expect the vertex with the maximal norm to \enquote{win}, but another
vertex with a smaller norm but larger cosine could also win. Suppose we make the angle
\( \theta \) larger; both cosines will increase, but the smaller cosine will increase
relatively more.\footnote{
	Suppose \( -\pi/2 < \alpha < \beta < -\varepsilon \); we claim that
	\begin{equation}
		\frac{\cos (\alpha + \varepsilon)}{\cos \alpha} >
		\frac{\cos (\beta + \varepsilon)}{\cos \beta}.
	\end{equation}
	We know that the derivative of \( \cos \) is decreasing, so
	\begin{equation}
		\frac{\cos (\alpha + \varepsilon) - \cos \alpha}{\varepsilon}
		> \frac{\cos (\beta + \varepsilon) - \cos \beta}{\varepsilon}   .
	\end{equation}
	We can use this in
	\begin{align}
		\frac{\cos(\alpha + \varepsilon)}{\cos \alpha} - 1
		= \frac{\cos(\alpha + \varepsilon) - \cos \alpha}{\cos \alpha} >
		\frac{\cos(\beta + \varepsilon) - \cos \beta}{\cos \alpha} >
		\frac{\cos(\beta + \varepsilon) - \cos \beta}{\cos \beta}=
		\frac{\cos(\beta + \varepsilon)}{\cos \beta} - 1.
	\end{align}
}
Therefore, the larger norms only gain more of an advantage as \( \theta \) increases, so
\begin{equation}
	\theta \mapsto \norm{\Sj 1}, \quad \theta \in \cc{0, \varepsilon}
\end{equation}
is increasing. With the analogous claim for the second walk,
\begin{equation}
	\theta \mapsto \norm{\Sj 2} - \norm{\Sj 1}, \quad
	\theta \in \cc{0, \varepsilon}
\end{equation}
is maximal for \( \theta=0 \).

To summarise this section:
\begin{itemize}
	\item on \( \cc{0, \varepsilon} \), we use the above;
	\item elsewhere, the arguments are almost the same as in the previous section.
\end{itemize}

\end{document}
